\documentclass[12pt,oneside]{book}

% Packages for cyberpunk noir styling
\usepackage[utf8]{inputenc}
\usepackage[T1]{fontenc}
\usepackage{geometry}
\usepackage{fancyhdr}
\usepackage{titlesec}
\usepackage{color}
\usepackage{xcolor}
\usepackage{graphicx}
\usepackage{tcolorbox}
\usepackage{enumitem}
\usepackage{hyperref}
\usepackage{fancyvrb}

% Cyberpunk color scheme
\definecolor{neonblue}{RGB}{0,255,255}
\definecolor{neongreen}{RGB}{57,255,20}
\definecolor{darkbg}{RGB}{20,20,20}
\definecolor{matrixgreen}{RGB}{0,200,0}
\definecolor{terminalgreen}{RGB}{0,255,0}

% Page geometry
\geometry{letterpaper,margin=1in}

% Headers and footers
\pagestyle{fancy}
\fancyhf{}
\fancyhead[C]{\color{neonblue}Shadows in the Chain}
\fancyfoot[C]{\color{matrixgreen}\thepage}

% Chapter styling
\titleformat{\chapter}[display]
  {\normalfont\huge\bfseries\color{neongreen}}
  {\chaptertitlename\ \thechapter}
  {20pt}
  {\Huge\color{neonblue}}

% Section styling
\titleformat{\section}
  {\normalfont\Large\bfseries\color{neongreen}}
  {\thesection}
  {1em}
  {}

% Terminal text box styling
\newtcolorbox{terminalbox}{
  colback=black,
  colframe=terminalgreen,
  fontupper=\ttfamily\color{terminalgreen},
  arc=0pt,
  boxrule=1pt
}

% Hyperlink setup
\hypersetup{
    colorlinks=true,
    linkcolor=neonblue,
    filecolor=neongreen,
    urlcolor=neonblue,
    pdftitle={Shadows in the Chain - Chapters 3-4},
    pdfauthor={Q-NarwhalKnight Authors},
    pdfsubject={Cyberpunk Noir Novel},
    pdfkeywords={cyberpunk, noir, blockchain, cryptocurrency, thriller, quantum computing}
}

\begin{document}

\begin{titlepage}
\centering
\vspace*{2cm}

{\Huge\bfseries\color{neongreen} Shadows in the Chain \par}
\vspace{0.5cm}
{\LARGE\color{neonblue} Chapters 3-4 \par}
\vspace{1.5cm}

{\Large\color{neonblue} Q-NarwhalKnight Authors \par}
\vspace{2cm}

{\color{darkbg}
\begin{tcolorbox}[colback=darkbg,colframe=neonblue,arc=0pt,outer arc=0pt]
\color{neongreen}\textbf{Genre:} Cyberpunk, Noir, Thriller\\
\color{neongreen}\textbf{Setting:} Zurich, Singapore, Global\\
\color{neongreen}\textbf{Status:} Draft - Chapters 3-4\\
\color{neongreen}\textbf{Word Count:} 8,000 words
\end{tcolorbox}
}

\vfill
{\large\color{matrixgreen} Generated by shadowchain-writer CLI \par}
{\color{neonblue} \today \par}

\end{titlepage}

\newpage
\tableofcontents
\newpage

% CHAPTER 3
\chapter{Quantum Entanglement}

\section*{Scene 1: The Zurich Dead Drop}

The morning light over Lake Zurich cut through the autumn mist like a blade through silk. Elena Voss stood on the Quaibrücke, her hands gripping the cold stone railing, watching the swans glide across water that mirrored the sky. She'd chosen this spot because it was public, exposed, and—most importantly—under constant surveillance by three different intelligence agencies. Paradoxically, that made it safer than any safehouse.

Her phone buzzed. Not a call. A blockchain notification.

\begin{terminalbox}
QUANTUM\_KEY\_EXCHANGE\_INITIATED.\\
COORDINATES REQUIRE DUAL\_VERIFICATION.
\end{terminalbox}

Elena's pulse quickened. She pulled out a paper map—analog, untraceable—and cross-referenced the coordinates with a pre-memorized location. The Altstadt. Old Town. Specifically, a centuries-old pharmacy near the Grossmünster cathedral.

She walked. No hurry. Zurich demanded elegance even from spies.

The pharmacy smelled of lavender and old wood. Behind the counter, an elderly woman with silver hair looked up from a ledger. Elena ordered a specific herbal remedy—one that didn't exist. The woman's expression didn't change, but her eyes flicked to a shelf in the back corner.

Elena browsed. Her fingers found a small, leather-bound journal wedged between a jar of dried chamomile and a bottle of tincture. Inside the journal: a single USB drive, wrapped in anti-static foil.

No words were exchanged. Elena paid cash for a tin of throat lozenges and left.

\vspace{1em}

Back at her hotel—a nondescript Ibis near the Hauptbahnhof—Elena inserted the drive into an air-gapped laptop running Tails OS. The files decrypted using a quantum-resistant algorithm: Crystals-Kyber, the NIST standard for post-quantum key exchange.

One video file. One text document.

She opened the video first.

A woman appeared on screen. Mid-thirties. Japanese. Exhausted eyes behind wire-rimmed glasses. Lab coat. Behind her, whiteboards covered in equations—Schrödinger operators, quantum gate diagrams, error correction codes.

``My name is Dr. Yuki Tanaka,'' the woman said, her voice barely above a whisper. ``I used to work at CERN. Before that, NIST. I helped design Crystals-Kyber.'' A bitter laugh. ``Ironic, isn't it? I made the lock. Now I'm watching someone else forge the key.''

Elena leaned closer.

``You're looking for Nexus Veil's architecture. The Master Node. But that's not the real problem anymore.'' Tanaka's hands trembled as she adjusted her glasses. ``Six months ago, I was approached by an organization. They called themselves Kronos. They wanted my expertise in quantum key distribution—specifically, BB84 protocol implementation for blockchain consensus.''

BB84. Elena knew it. The first quantum cryptography protocol, invented in 1984 by Bennett and Brassard. Theoretically unbreakable because it relied on the laws of physics: measuring a quantum state collapses it, making eavesdropping detectable.

``They're not just building a secure blockchain,'' Tanaka continued. ``They're building a post-quantum financial system. One that can survive Q-Day—the day quantum computers break RSA and ECC encryption. But here's the catch: Kronos isn't waiting for Q-Day. They're accelerating it.''

The video cut to footage of a laboratory. Rows of cryogenic chambers. Inside: quantum processors. Dozens of them. Not the noisy, error-prone prototypes Elena had read about. These were stable. Scalable. Operational.

``They've achieved something we thought was a decade away: fault-tolerant quantum computing. Sixty-eight logical qubits. Enough to run Shor's algorithm. Enough to break RSA-2048 in hours instead of centuries.''

Elena's blood ran cold.

``If Kronos controls both the quantum computers and the post-quantum infrastructure, they control the transition. Every bank, every government, every individual will have to migrate to their system. And they'll have a backdoor into everything.''

The video ended with coordinates. Singapore. A research facility listed as a biotechnology lab but operating under shell corporations registered in the Caymans.

Elena opened the text file. It contained technical specifications: qubit architectures, error rates, estimated timelines for scaling. And a name: Dr. Rajesh Krishnan. Quantum physicist. Last known location: Singapore. Status: Missing.

\section*{Scene 2: The Zurich Protocol}

Elena met Dr. Tanaka in person three hours later, at a different location: a conference room in the University of Zurich's physics department. Tanaka had insisted on neutral ground, surrounded by academia and witnesses.

``You took a risk coming here,'' Elena said, sitting across from her at a long table covered in printouts and diagrams.

``I'm already dead if Kronos finds me,'' Tanaka replied. She looked even more exhausted in person. ``At least this way, someone knows.''

Elena gestured to the papers. ``Walk me through it. Quantum key distribution. BB84. Why does it matter for Nexus Veil?''

Tanaka pulled out a marker and approached a whiteboard. She drew two points labeled Alice and Bob.

``BB84 is simple in theory. Alice sends Bob a stream of photons, each encoding a bit of information. The photons are polarized—vertical, horizontal, diagonal. Alice randomly chooses polarization bases. Bob measures them using randomly chosen bases too.''

She drew arrows, circles, measurement axes.

``After transmission, Alice and Bob compare bases publicly—but not the results. Where their bases matched, the measurements are guaranteed to be identical. That's their shared secret key.''

``And if someone eavesdrops?'' Elena asked.

``That's the beauty. Quantum mechanics says measuring a photon collapses its state. If an eavesdropper intercepts the photon, they disturb it. Alice and Bob will detect anomalies in their key—bit errors above the noise threshold. They abort and try again.''

Elena frowned. ``So it's unbreakable.''

``In theory.'' Tanaka capped the marker. ``In practice, it's hardware-dependent. You need perfect single-photon sources, perfect detectors, perfect transmission. Any imperfection creates side channels. And then there's the authentication problem.''

``Authentication?''

``BB84 doesn't authenticate the parties. You need a pre-shared secret to verify that Alice is really Alice and Bob is really Bob. Otherwise, a man-in-the-middle attack is trivial.'' Tanaka sat back down. ``That's where Nexus Veil comes in. It uses blockchain consensus to distribute authentication keys. Decentralized, post-quantum secure, theoretically tamper-proof.''

``But Kronos has a backdoor,'' Elena said.

``Not a backdoor. A master key. The genesis block of Nexus Veil contains a Shamir secret-sharing scheme—a cryptographic puzzle split across multiple holders. Reassemble enough shards, and you can rewrite the consensus rules. Mint new tokens. Revoke authentication. Control the entire network.''

``That's the Master Node,'' Elena murmured.

``Exactly. And Kronos has been collecting the shards. They have three out of five. One more, and they achieve majority control.''

Elena's mind raced. ``Who holds the other two?''

``One is in Singapore. Krishnan has it. The other...'' Tanaka hesitated. ``Marcus Hale. CIA.''

Elena's stomach dropped.

Marcus.

Her nemesis. Her mirror image. The man who'd burned her twice and nearly killed her once.

``Why would the CIA have a shard?''

``Because they helped build Nexus Veil,'' Tanaka said quietly. ``Not officially. But they seeded it. Funded it. Used it to track dark money, terror financing, sanctions evasion. They thought they could control it. They were wrong.''

\section*{Scene 3: The Kronos Intercept}

Elena left the university through a maintenance exit, her mind spinning. Marcus Hale. Of course he was involved. The man collected secrets like other people collected stamps.

But if Kronos knew about the shards, they'd be hunting Marcus too. And if Marcus was in play, that meant the CIA was compromised or desperate or both.

She took a tram to the Hauptbahnhof, blending into the evening commuter crowd. She needed to think. Plan. Decide whether to trust Tanaka's intel or—

Her instincts screamed a fraction of a second before it happened.

The man was nondescript. Beige jacket, jeans, wire-rimmed glasses. He stepped onto the tram at the next stop, holding a newspaper. He didn't look at her.

But Elena knew.

She moved three stops later, exiting at Paradeplatz. The financial district. Crowds, cameras, security. She walked briskly toward the Bahnhofstrasse, Zurich's main shopping street.

The man followed.

Elena ducked into a Lindt chocolate shop, weaving through tourists. Out the side exit. Down an alley. She paused at the corner, listening.

Footsteps. Multiple sets.

Kronos wasn't subtle this time. They were herding her.

She ran.

The alley opened onto a narrow street paralleling the Limmat River. Elena sprinted, her breath fogging in the cold air. Behind her: shouts, footsteps, the unmistakable sound of a suppressed pistol.

A bullet sparked off the cobblestones inches from her foot.

Elena vaulted a railing, landing hard on the riverbank below. Her ankle screamed in protest, but she kept moving. The river path was deserted—no witnesses, no cameras, no help.

Another shot. This one closer.

She spotted a maintenance ladder leading up to a bridge. She climbed, her injured ankle threatening to give out. At the top, she pulled herself onto the roadway and flagged down a taxi.

``Flughafen,'' she gasped. ``Airport. Now.''

The driver hesitated, seeing her disheveled state.

Elena slapped a hundred-franc note on the dashboard. ``Now.''

The taxi sped off.

\section*{Scene 4: The Singapore Gambit}

Elena didn't go to the airport. She had the driver drop her at a parking garage in Oerlikon, where she kept a second go-bag in a rented locker. New passport. New clothes. New identity.

She took a train to Munich, then a flight to Istanbul, then another to Singapore. Three countries, three identities. Kronos would track her eventually, but she'd bought time.

Singapore was heat and humidity and the smell of durian. Elena checked into a capsule hotel in Little India, paying cash. She slept for four hours, her hand on a knife under the pillow.

When she woke, she pulled out a burner phone and composed a message to an encrypted drop box—one she knew Marcus Hale monitored.

\begin{terminalbox}
Kronos has three shards. They're coming for yours.\\
Singapore. Krishnan. 48 hours.\\
--Ghost
\end{terminalbox}

She hesitated before hitting send. Reaching out to Marcus was insane. He'd betrayed her in Prague. Nearly killed her in Istanbul. He'd choose the mission over her every single time.

But he was also the only person who understood the stakes.

She sent the message.

Then she went hunting for Dr. Rajesh Krishnan.

\newpage
% CHAPTER 4
\chapter{The Singapore Protocol}

\section*{Scene 1: The Newton Circuit}

Singapore's Newton Food Centre was chaos rendered edible. Steam rose from fifty different hawker stalls, mingling with the clatter of metal spatulas, the hiss of woks, and the polyglot chatter of a dozen languages. Elena sat at a plastic table under harsh fluorescent lights, nursing a cup of teh tarik and watching the crowd.

Dr. Rajesh Krishnan was supposed to meet her here. Tanaka had provided the contact protocol: a specific stall (Hainanese chicken rice, stall 37), a specific time (7 PM), a specific phrase (``Is the rice fresh today?'').

Elena had arrived early. Scouted exits. Identified security cameras. Marked potential threats.

Krishnan was late.

At 7:15, a thin man in his forties approached stall 37. Wire-rimmed glasses. Nervous energy. A leather satchel slung over one shoulder. He ordered food, glancing around too obviously.

Elena waited. Watched.

No tail. Or at least, no obvious tail.

She stood and walked to the stall, positioning herself behind Krishnan in line.

``Is the rice fresh today?'' she asked.

Krishnan flinched. Turned. His eyes widened behind his glasses. ``You're... Tanaka sent you?''

``Keep your voice down,'' Elena murmured. ``Order your food. Sit at table twelve. Don't look at me.''

He nodded jerkily.

Five minutes later, they sat across from each other, separated by plates of chicken rice and chili sauce. Krishnan picked at his food, not eating.

``Tanaka said you'd help me,'' he whispered. ``She said... she said you could get me out.''

``Depends on what you're carrying,'' Elena said.

Krishnan's hand went instinctively to his satchel. ``The shard. The quantum key fragment. I've kept it offline, encrypted, distributed across three flash drives using a threshold scheme. Even if Kronos intercepts one or two, they can't reconstruct it.''

``Smart. Why haven't you gone to the authorities?''

``Because I don't know who to trust.'' Krishnan's voice cracked. ``Kronos has people everywhere. Government. Academia. Intelligence services. They approached me a year ago, offering funding for my research—quantum error correction for distributed systems. I thought it was legitimate. By the time I realized what they were building, it was too late.''

``What are they building, exactly?''

``A quantum supremacy network.'' Krishnan leaned forward. ``Not just a single quantum computer. A distributed system of quantum processors linked via entanglement-based communication. Imagine blockchain consensus, but powered by quantum mechanics. Instant. Unhackable. Scalable to millions of nodes.''

Elena frowned. ``Quantum entanglement can't transmit information faster than light. No-communication theorem.''

``Correct. But it can establish correlations instantly. Use those correlations for cryptographic protocols—quantum key distribution, Byzantine fault tolerance, consensus mechanisms. Kronos isn't trying to break the laws of physics. They're bending them.''

``And the timeline?''

``Six months until they achieve stable 4096-qubit systems. At that scale, RSA-2048 is breakable in hours. Elliptic curve cryptography becomes worthless. Bitcoin, Ethereum, every blockchain—vulnerable.'' Krishnan's hands trembled. ``And then they flip the switch. Deploy Nexus Veil globally as the 'solution.' Every financial institution, every government migrates to their system. And Kronos controls the Master Node.''

Elena processed this. ``How do we stop them?''

``You can't stop them. But you can delay them. Disrupt their supply chains—quantum processors require rare-earth elements, isotopically pure silicon, dilution refrigerators. Sabotage their labs. Leak their research to force transparency.'' He paused. ``Or use the Master Node against them. Rewrite the consensus rules before they do. Burn the whole thing down.''

\section*{Scene 2: The Quantum Threat}

They left the hawker center separately. Krishnan gave Elena a key to a storage locker at the Bugis MRT station—one of the flash drives was there. The other two were in locations he'd reveal only after confirming Elena's extraction plan.

She took the Circle Line to Bugis, riding three stops past her destination before doubling back. No tail. Not yet.

The locker was in a dimly lit row near the public restrooms. Elena retrieved the flash drive—a generic 32GB USB stick, indistinguishable from a thousand others. She pocketed it and left.

Her burner phone buzzed. A message from her encrypted drop box.

\begin{terminalbox}
Singapore. 48 hours. Rendezvous at the Fullerton.\\
Come alone. We need to talk.\\
--Hale
\end{terminalbox}

Elena stared at the screen. Marcus was here. In Singapore. And he wanted to meet at the Fullerton Hotel—a five-star property on the waterfront, crawling with security and surveillance.

Neutral ground? Or a trap?

\section*{Scene 3: The Fullerton Gambit}

The Fullerton was all marble and brass and colonial-era grandeur. Elena entered through the main lobby, dressed in business casual—blazer, slacks, minimal jewelry. She looked like a consultant or lawyer, forgettable in a sea of forgettable professionals.

Marcus had specified the bar. 9 PM.

Elena arrived at 8:45. She ordered sparkling water and positioned herself where she could watch both entrances and the elevators.

Marcus appeared at 9:03. Taller than she remembered. Gray creeping into his temples. The same sharp eyes, the same predatory calm.

He sat across from her without a word. The bartender brought him a whiskey, neat.

``You look terrible,'' Marcus said.

``You look old,'' Elena replied.

His lips twitched. Almost a smile. ``Zurich was messy. Kronos doesn't usually deploy wet teams that openly.''

``They're desperate. Three shards down, two to go.''

``Three and a half,'' Marcus corrected. ``They have partial access to one of Krishnan's fragments. He doesn't know it yet, but his air-gapped system was compromised via a supply chain attack. Infected USB controller firmware. Slow exfiltration over six months.''

Elena's stomach dropped. ``Then Krishnan is—''

``Exposed. Yes. And by meeting him, so are you.'' Marcus sipped his whiskey. ``Kronos knows you're here. They're letting you run because you're leading them to the other shards.''

``Why warn me?''

``Because I need something from you.'' Marcus leaned forward. ``Krishnan has information about Kronos's leadership structure. Specifically, the identities of the Architects—the core group funding and directing the quantum supremacy project. If we take them out, we disrupt the timeline. Buy the world a few years.''

``And your shard?''

``Leverage. I'm not handing it over to Kronos, and I'm not handing it to you. But I will use it as bait.'' He slid a phone across the table. ``There's a location on here. A facility in Moscow. Cold War-era bunker, repurposed as a quantum research lab. One of the Architects operates out of there. We hit it together, extract the intel, and I'll consider sharing my shard.''

Elena studied him. ``You're lying.''

``Probably.'' Marcus smiled. ``But you don't have a better option. Krishnan's compromised. Tanaka's on the run. And Kronos is twelve moves ahead. So you can work with me, or you can die alone in a Singapore alley.''

Elena picked up the phone. The location was encrypted, locked behind a biometric scan. Marcus's fingerprint or retina, she guessed.

``When do we leave?'' she asked.

``Forty-eight hours. I'll send you the details.'' Marcus stood. ``And Elena? Don't try to double-cross me. You won't like how it ends.''

He left.

Elena sat alone, staring at the phone.

\section*{Scene 4: The Hawker Center Trap}

Elena returned to the Newton Food Centre the next evening to meet Krishnan again—he'd promised the locations of the other two flash drives. But as she approached table twelve, her instincts screamed.

Krishnan was there. But his body language was wrong. Stiff. Afraid.

And sitting two tables away, partially obscured by a pillar: a woman in a white blouse, reading a newspaper. Asian features. Sharp posture. Expensive watch.

Elena's mind cataloged the details in a fraction of a second. The woman wasn't eating. Wasn't sweating despite the heat. And her newspaper was in English—an outlier in a crowd dominated by Mandarin and Tamil.

MSS. Chinese Ministry of State Security.

Elena walked past table twelve without stopping. Krishnan's eyes followed her, pleading. But she couldn't help him. Not yet.

She circled the hawker center, exiting through the back near the restrooms. The woman from table twelve didn't follow. Someone else did—a young man in a polo shirt, carrying a takeout bag.

Elena led him into a narrow alley behind the kitchens. Garbage bins. Steam vents. No cameras.

The man reached for something under his shirt.

Elena moved first. She grabbed his wrist, torqued it sideways, drove her knee into his ribs. He gasped, doubling over. She wrenched the pistol from his waistband and slammed him against the wall.

``Who sent you?'' she hissed.

He spat blood. Said nothing.

Elena pressed the muzzle of the pistol against his kneecap. ``Last chance.''

``Kronos,'' he gasped. ``They... they have Krishnan's family. Wife. Two kids. He cooperates, or they die.''

Elena's heart sank.

``Where?''

``I don't know. I'm just surveillance. Please—''

Elena pistol-whipped him. He collapsed.

She returned to the hawker center cautiously. Krishnan was gone. So was the woman.

Elena cursed under her breath. Krishnan was compromised. The flash drive in her pocket might be the only shard fragment not under Kronos's control—unless Marcus was right, and even that was tainted.

She pulled out her burner phone and composed a message to Marcus.

\begin{terminalbox}
Krishnan burned. Family hostage. Moscow is our only play.\\
--Ghost
\end{terminalbox}

She hit send.

Then she disappeared into the Singapore night.

\section*{Scene 5: The Moscow Decision}

Elena spent the next thirty-six hours moving. She switched hotels three times. Burned two identities. Acquired a forged Russian visa through a contact in Jakarta.

Marcus sent the Moscow details via encrypted message: coordinates, entry protocols, a contact name. The facility was beneath Gorky Park, accessible through a decommissioned metro tunnel. Heavily guarded. High risk.

Possibly a suicide mission.

Elena sat in a cramped hotel room in Geylang, staring at the message on her laptop screen. Marcus's plan was insane. But so was trusting Krishnan, and that had already blown up. So was trusting Tanaka, who might be Kronos herself for all Elena knew.

Every option was a trap.

But doing nothing meant Kronos won.

Elena closed the laptop. She packed her go-bag—passport, cash, weapons, the flash drive. She booked a flight to Moscow via Dubai, departing in six hours.

Before she left, she sent one final message to Tanaka.

\begin{terminalbox}
If I don't make it back, the shard is yours. Use it wisely.\\
Coordinates attached.\\
--Ghost
\end{terminalbox}

She didn't believe Tanaka would use it wisely. But at least it wouldn't fall into Kronos's hands by default.

Elena shouldered her bag and walked out into the humid Singapore night. She had forty-eight hours to reach Moscow, infiltrate a quantum research facility, and extract intel on the Architects.

And after that?

She'd improvise.

She always did.

\newpage
\appendix
\chapter{Technical Notes}

\section*{Quantum Key Distribution (BB84 Protocol)}

BB84 is the first quantum cryptography protocol, proposed by Charles Bennett and Gilles Brassard in 1984. It uses the principles of quantum mechanics to establish a shared secret key between two parties (Alice and Bob) that is provably secure against eavesdropping.

\textbf{Key Properties:}
\begin{itemize}
\item Uses quantum states (polarized photons) to encode information
\item Eavesdropping is detectable due to quantum measurement disturbance
\item Requires authenticated classical channel for basis reconciliation
\item Theoretically information-theoretically secure
\end{itemize}

\section*{Post-Quantum Cryptography}

\textbf{Crystals-Kyber:} NIST-standardized lattice-based key encapsulation mechanism (KEM) designed to resist quantum computer attacks.

\textbf{Crystals-Dilithium:} NIST-standardized lattice-based digital signature scheme for post-quantum security.

\textbf{Timeline:} NIST finalized post-quantum cryptography standards in 2024. Migration timelines vary by industry (5-10 years for full deployment).

\section*{Quantum Computing Threat Model}

\textbf{Shor's Algorithm:} Quantum algorithm that can factor large integers and solve discrete logarithms in polynomial time. Breaks RSA, DSA, ECDSA.

\textbf{Grover's Algorithm:} Quantum algorithm that provides quadratic speedup for unstructured search. Weakens symmetric encryption (AES-128 → AES-256 recommended).

\textbf{Q-Day:} Hypothetical date when quantum computers become powerful enough to break current cryptographic standards (estimated 2030-2040, varies by source).

\end{document}
